% abstract.tex
\begin{abstract}
Existing audio-to-MIDI tools recover discrete note events from recordings but discard the timbral characteristics that define an instrument's identity.
We present \textsc{Instrumental}, a system that recovers continuous synthesizer parameters from audio by coupling a differentiable 28-parameter subtractive synthesizer with CMA-ES, a derivative-free evolutionary optimizer.
The synthesizer implements a full subtractive signal chain -- four oscillator types with unison voicing, a resonant low-pass filter with dedicated ADSR, two-band parametric EQ, and reverb -- entirely in PyTorch.
We optimize a composite perceptual loss combining mel-scaled multi-resolution STFT, spectral centroid distance, and MFCC divergence, achieving a matching loss of 2.09 on real recorded audio.
We systematically evaluate eight hypotheses for improving convergence -- including Chebyshev waveshaping, FM synthesis, alternative perceptual losses, and multi-start strategies -- and find that only parametric EQ boosting yields meaningful improvement.
Our results demonstrate that CMA-ES outperforms gradient descent on this non-convex landscape, that more parameters do not monotonically improve matching, and that spectral-analysis initialization accelerates convergence over random starts.
\end{abstract}
